\chapter{The Corridor Rotates}

\section{Rotation turns number into phase}

The last chapter named the number: a box, a discrete label the pigeonhole forced. This chapter shows that box
beginning to \emph{rotate}. Given a second axis to turn in, the number stops being a mark on a line and becomes a
\emph{phase} --- a position on a circle. And the construction, it turns out, has been carrying the second axis all
along, unused. This chapter opens by putting it to work; and in doing so it finally grounds the one reading the book
has so far only named --- the sign it read on the electron's box.

Start with what the construction actually carries, because the second axis is not new machinery. There is a type the
book has leaned on throughout --- the type of a \emph{real} axis --- and at the higher structures the construction
binds it not once but \emph{twice}: once under the name \emph{real}, and again, beside it, under the name
\emph{imaginary}. Two instances of the same axis type, standing side by side. That is a plane, not a line. The
imaginary axis was not smuggled in when it was needed; it was there the whole time, a second copy of the real that the
corridor carried without using. This is the whole of the derived content, and it is thin and honest: there are two
axes, and a number that has lived on one of them can be given a coordinate on the other.

Now the reading, and it must be marked as exactly that. To carry the number from the real axis onto the imaginary one
is a quarter-turn, and the generator of that quarter-turn is written $-i$. Relabelling which axis the number sits on
--- turning \emph{real} into \emph{imaginary} --- lifts it off the straight line and sets it on a circle, and a
position on a circle is a \emph{phase}: an angle, an orientation, rather than a magnitude read off a ruler. The
discrete box the last chapter forced becomes, on the plane, an angle. But here the book owes the reader its sharpest
line of honesty, and it is the line that governs the whole chapter. \emph{The machine builds the plane; the book
supplies the turn.} What the code carries is two static axes --- two instances of the real, sitting there. It carries
no rotation: no $-i$, no quarter-turn, no circle, no exponential of an angle, nowhere in it. The plane is derived; the
\emph{turn} that carries one axis into the other is the reading, laid over the two-axis structure by the book. The
imaginary unit and its quarter-turn are a mathematical bridge --- a true and standard one --- but a bridge the book
builds, not a claim the machine proves.

And now the reading pays a debt. Two chapters ago the book read the electron's box as carrying a sign --- minus one ---
and marked the sign as its own reading, promising the grounding for later. Here is the grounding. Two quarter-turns,
taken in a row, make a half-turn; and a half-turn on the circle is a reversal --- it sends a point to its opposite. So
the minus the book read on the box is exactly a \emph{doubled quarter-turn}: the corridor rotated twice, come round to
the far side, which is negation. The sign is the loop closing on the plane rather than on the line. And the discipline
holds firm across the grounding: the box itself --- the discrete, forced integer the pigeonhole proved --- stays
proved, untouched; it is the \emph{sign's} reading, the minus as a doubled turn, that is the interpretive bridge, and
it stays marked. The full family of quarter-turns around the circle --- the four positions the rotation visits --- is
the spinor, and the spinor is physics; it belongs to the next volume. This chapter grounds the sign in a rotation and
hands the spin forward.

In the two verbs, rotation is a relabelling seen from a new side. To carry a number from the real axis to the imaginary
is to re-select which characteristic --- which axis --- names it: a tange performed not to tell one number from another
but to turn one number around. The count did not change; its label rotated. And because a rotation preserves length,
the magnitude the number carried is untouched --- only its orientation, its place on the circle, changed. So this is
where the two verbs the book has run on become \emph{orientation}: the phase is the number re-tanged onto a second
axis, the same object relabelled around a circle, its size exactly what it was. Telling apart, which is where the book
began, has learned to turn.

So the corridor begins to rotate. The number that was a discrete box is now a box with an angle available to it, and
the sign it carried has a home: a doubled quarter-turn on a plane the machine built and the book learned to turn in.
What the volume proves stays what it proved --- the box, the forced integer --- and what it reads, it marks --- the
turn, the phase, the sign as a rotation. The rest of the chapter follows the turn: what it costs to relabel, and what,
on the far side of the rotation, the machine finds it has been measuring.

\section{What a decoding is depends on the frame}

The last section rotated the frame and turned a number into a phase. This section draws the lesson the rotation
teaches, and it is a lesson about where a decoding's character lives. The same residue is, we have seen, a label, a
count, an order, a parity, an orientation --- and the point of this section is that which one it \emph{is} is not
settled by the residue. It is settled by the frame that reads it. And the sharpest form of the claim is the one the
source states outright: whether a decoding runs covariant or contravariant --- the two bands the faces cut along ---
is \emph{not a property in the object at all}. It is the object's \emph{response} to a change of frame.

The source puts it as plainly as it can, and because it says it in its own words the book can hand the principle over
as stated rather than argued. A specialist understands the difference between two kinds of quantity, it says, because
she has watched them under a stimulus: apply a change of coordinates, and one kind goes \emph{with} the change while
the other goes \emph{against} it. The two are not intrinsically different objects, waiting to be told apart by
inspection. They differ only in how they \emph{respond} when the frame is changed. And then the line the whole section
turns on: \emph{the covariance is not in the definition; the covariance is in the response to the stimulus.}
``Covariant'' and ``contravariant'' are not labels stamped on a thing. They are names for how a thing behaves when you
change how you look at it.

And the same principle wears three faces, one for each specialist who meets it. To the geometer the stimulus is a
change of coordinates, and the response is whether a quantity goes with the change or against it. To the category
theorist the stimulus is a morphism, and the response is whether a construction preserves its direction or reverses it.
To the type theorist the stimulus is a subtyping relation, and the response is whether a constructor respects the
ordering or flips it. Three specialists, the source says; three stimuli; three response patterns; one word. It is one
variance concept throughout, and which of its patterns you see is fixed entirely by which stimulus --- which frame ---
you bring. Change the frame, and the same object shows you a different face. These three frames are all mathematical
--- coordinates, morphisms, subtyping --- and the book keeps to them; the frames that read the residue as a physical
quantity belong to the next volume.

Now cash the principle against the residue. Its decodings --- label, count, order, parity, and the orientation the last
section added --- are frame-relative, every one. The same object \emph{is} a count under one stimulus, an order under
another, a phase under a rotation; none of these is more truly the object than the rest, because none of them is in the
object alone. Each is the object's response to a particular frame. And this reaches back and explains the faces of an
earlier chapter. The covariant-or-contravariant band a reading falls on is its response to the fact-parity stimulus;
the four faces are four \emph{frames}, four different stimuli brought to one residue. So what the residue \emph{is}
depends on which gauge asks --- and this is exactly why the agreement of the faces is evidence and not luck. Four
frames apply four stimuli to one object, and the object responds \emph{consistently} across all of them. A thing that
answered four independent stimuli the same way by accident would be a miracle; a single object, responding as itself
under each frame, is the ordinary explanation --- and it is the one the agreement points to. (Even the calibration
frame the book will finally settle on is a choice of stimulus, not a neutral vantage; there is no frame-free place to
read from, and the book does not pretend there is one.)

In the two verbs the lesson is sharp, and it revises what the verbs are. Covariant is a \emph{funge}: the decoding goes
with the frame, counting with the world's parity. Contravariant is a \emph{tange}: it goes against, our label laid over
the object and inverted relative to it. But --- and this is the turn --- \emph{which one a decoding is depends on the
frame.} The same object funges under one stimulus and tanges under another; the verb it earns is not a fixed attribute
it carries but its response to how it is read. So this is where the two verbs the book has run on stop being properties
of the residue and become \emph{verdicts relative to a frame} --- the same object, funged here and tanged there,
depending only on which stimulus asks. Telling apart and bagging by likeness are not things the object is; they are
things it \emph{does}, when a frame acts on it.

So the rotation of the last section generalizes into a rule for the whole instrument: a reading is never frame-free,
and what a decoding \emph{is} is always the answer to a \emph{how are you looking}. The book will lean on this at the
end twice over --- once to explain why four faces agreeing is proof, and once to admit that even its final, careful
frame is a chosen one. The next section asks what the choosing costs.

\section{Relabelling has a cost}

Two sections have rotated the frame and made the reading depend on it. This closing section adds the fact that keeps
the rotation honest, and it is a plain one: relabelling is \emph{not free}. To read the number in a rotated frame, the
machine cannot simply declare the new name; it has to \emph{re-elaborate} the number --- compute it again, in the new
frame --- and computation is billed. The turn of the first section costs the machine time. And so this chapter, which
began by turning a number into a phase, ends by pricing the turn.

The cost is not a metaphor; the machine literally measures it. The construction carries a facility that meters, for any
term, the \emph{elaboration cost} of that term --- the work the checker does to reduce it to a settled form, counted in
the machine's internal beats. To relabel a number is to re-elaborate it under a new frame, and re-elaboration runs that
meter. The rotation, then, has a price the machine actually pays: so many beats to re-read the number on the second
axis. A change of name that looked, in the first section, like a free quarter-turn is, underneath, a computation with a
bill attached. (This meter is the seed of a later chapter's central turn --- there the machine turns the meter on
itself; here it only establishes that relabelling costs.)

And the currency the bill is written in is worth naming, because it is not an external fee. The meter counts
\emph{heartbeats} --- the machine's own native tick, the same constant beat that, chapters ago, repeatability turned
into time. So relabelling is paid for in exactly the currency the machine runs on: its own clock. The phase of the
first section does not cost some outside currency; it costs the machine a measurable number of its own beats to
compute, and the count is calibrated against the machine's own fixed constant, not against anything outside it. This is
a small fact with a long reach. Because every act the machine performs is paid for in its own clock, the machine can
later be turned to meter \emph{itself} --- and the deepest number the whole instrument reaches will be read not as a
magnitude out in the world but as a \emph{cost}, a price paid in beats, the machine weighing what it spends to know its
own account. Every act, even a mere relabelling, is billed in the machine's own time.

There is a physical name for what a continuous relabelling costs, and the book names it once and hands it on. A
symmetry you can perform by degrees --- a rotation you can turn a little at a time --- carries, by a theorem of
Noether's~\cite{noether1918}, a \emph{conserved current}: a quantity that the turning leaves invariant. The physical
gauge has a name for the conserved current of a rotation, and the name is \emph{angular momentum}. But that is physics,
and it is the next volume's to read. This volume keeps the part it can prove --- that the relabelling is billed, in
beats, in the machine's own currency --- and hands the conservation, and its physical name, forward. What this volume
stakes is a cost; what the conserved quantity \emph{is}, in the world, is the next gauge's.

In the two verbs the chapter closes with a price on each. The relabelling this chapter has studied --- the tange of a
number onto a second axis, its response to a new frame --- is billed: every re-selection under a new frame spends the
machine's beats. The tange is not free; to pick a thing out afresh, in a turned frame, is to pay for the picking. And
the funge has its own place in the accounting: what the continuous relabelling \emph{preserves} --- the invariant the
turning cannot change --- is exactly what survives being bagged across the rotation, the characteristic the funge keeps
no matter how far the frame is turned. So the chapter of the verbs-as-orientation ends with a ledger: orientation
costs, and conservation is the name for what remains when the bill is paid.

So the corridor, having learned to rotate, learns that rotating is not free. The number can be turned onto its second
axis and read as a phase --- but the turn is a computation, and the computation is billed in the machine's own beats.
What the volume proves is the cost; what the cost \emph{is}, physically --- a conserved current, an angular momentum
--- it marks and hands to the next gauge. And the small mechanical fact planted here, that the machine pays for its own
acts in its own clock, is the one the book will need at its climax, when the machine turns that same meter on itself
and reads its deepest number as the price of knowing what it is.
